% Nota independiente; texto revisado directamente en este archivo.
% Compilar desde este directorio con:
%   latexmk -pdf calculos_para_rogelio.tex
% No requiere archivos del manuscrito ni una bibliografia externa.
\documentclass[11pt,a4paper]{article}
\usepackage[T1]{fontenc}
\usepackage[utf8]{inputenc}
\usepackage[spanish,es-nodecimaldot,es-noshorthands]{babel}
\usepackage{lmodern,microtype,textcomp}
\usepackage[margin=24mm]{geometry}
\usepackage{amsmath,amssymb}
\usepackage{booktabs,longtable,array,calc}
\usepackage{fancyvrb,fvextra}
\usepackage{xcolor}
\usepackage{xurl}
\usepackage[colorlinks=true,linkcolor=blue!45!black,urlcolor=blue!45!black]{hyperref}
\DeclareUnicodeCharacter{00B5}{\textmu}
\DeclareUnicodeCharacter{00B0}{\textdegree}
\DefineVerbatimEnvironment{verbatim}{Verbatim}{breaklines=true,breakanywhere=true,fontsize=\small}
\providecommand{\tightlist}{\setlength{\itemsep}{0pt}\setlength{\parskip}{0pt}}
\setlength{\parindent}{0pt}
\setlength{\parskip}{5pt plus 1pt}
\setlength{\emergencystretch}{3em}
\setlength{\LTpre}{8pt}
\setlength{\LTpost}{8pt}
\title{Respuesta y ruido del CSA para NAUSICA}
\author{J. J. G\'omez-Cadenas}
\date{\today}
\begin{document}
\maketitle

{\bf Astra:} Los cálculos se han realizado con apoyo de Codex, versión Astra, high.

{\bf Idea:} El objetivo es determinar cuántos iones deben llegar a un pad durante
10 ms para medir su carga con una relación señal/ruido mayor que un cierto valor (por ejemplo, $S/N>5$).
Para ello el cálculo aplica al circuito LTspice una corriente calculada a partir del movimiento
de los iones. Después se filtra el voltaje de salida y se calcula el ruido
que queda con ese mismo filtro.

{\bf Spice}: Se usa LTspice 17.2.4 nativo en macOS, a 20 °C.


\section{Circuito simulado}

El cálculo parte del esquema \texttt{CSAfolded180nm.asc} y de los modelos de
\texttt{180nm\_professional.inc}. Se conservan las dimensiones y polarizaciones
de los transistores y la capacitancia de realimentación de 3 fF. \emph{El cálculo cambia tanto la fuente de corriente como
la capacitancia externa \texttt{Cd1}.}

Para los pads de 0.5 mm se adopta una capacitancia externa de 1.25 pF,
estimada por escalado de área a partir de Topmetal-S\footnote{Topmetal-S da 5 pF de capacidad para pads de 1 mm. Me limito a escalar por el área ---la nuestra es 0.5 mm---. }. Las capacitancias de
los transistores se suman a esta carga externa.

\section{Corriente de entrada}

La corriente se obtiene a partir de los iones que llegan al pad. El cálculo
está documentado en el artículo sobre NAUSICA. Para un ion de carga
$+e$, el potencial de ponderación $\phi_w$ da
\[
 i_1(t)=e\frac{d\phi_w}{dt},\qquad
 Q_g=\int i_1(t)\,dt
 =e\,[\phi_w(\mathrm{final})-\phi_w(\mathrm{inicio})].
\]
Para $N$ llegadas uniformes durante $\Delta=10\,\mathrm{ms}$, se usa
\[
 i_N(t)=\frac{N}{\Delta}\int_0^\Delta i_1(t-t_a)\,dt_a,
\]
donde $t_a$ es el instante de llegada al pad. Esta corriente se exporta
como una tabla de tiempos y valores para una fuente PWL de LTspice.

\section{Medida de la señal}

La carga llega durante un intervalo de 10 ms (el tiempo típico de un "bucket" correspondiente a 1 mm, se puede extender a 20 mm sin problema).

Se mide el cambio de voltaje mediante promedios de duración $T$ centrados
en sus extremos:
\[
 \overline V(t_n)=\frac{1}{T}
 \int_{t_n-T/2}^{t_n+T/2}V(t)\,dt.
\]
Se forma la diferencia corregida por descarga
\[
 D_n=-\left[\overline V(t_n)-\alpha\overline V(t_{n-1})\right],
 \qquad \alpha=\exp(-\Delta/\tau_0),
 \qquad \tau_0=0.3\,\mathrm{s}.
\]
El factor $\alpha$ descuenta la señal anterior que habría sobrevivido a la
descarga en el CSA ideal. Los 0.3 s son el valor elegido para esta corrección;
la descarga del circuito LTspice la determina el transistor de realimentación.
El signo hace positiva la señal.

Se calibra la medida con el propio circuito: $g$ es el valor de $D_n$ para
una corriente uniforme normalizada a un ion. En el régimen lineal, $N$ iones
producen $Ng$, y se estima su número como $\widehat N=D_n/g$.

Se toma $T=300\,\mu\mathrm{s}$ como duración de cada promedio. Implementar la medida
requiere leer el voltaje durante esas ventanas, que se extienden $T/2$
a cada lado del extremo del frame.

\section{Del espectro de ruido al ENC}

LTspice proporciona la densidad espectral de amplitud de ruido de salida
$a_V(f)$, en $\mathrm{V}/\sqrt{\mathrm{Hz}}$. La varianza después del
filtrado es
\[
 \sigma_D^2=\int_{0.1\,\mathrm{Hz}}^{1\,\mathrm{MHz}}
 a_V^2(f)\,|H(f)|^2\,df.
\]
Para la lectura de frames,
\[
 |H(f)|=\left|\operatorname{sinc}(fT)
 \left[1-\alpha e^{-2\pi i f\Delta}\right]\right|,
 \qquad \operatorname{sinc}(x)=\frac{\sin(\pi x)}{\pi x}.
\]
El primer factor corresponde al promedio y el segundo a la diferencia entre
lecturas. La expresión incluye la correlación del ruido entre ambas ventanas.
Se usa el ruido total \texttt{V(onoise)} calculado en el punto de trabajo
sin señal.

Al dividir el ruido de la medida por la señal de un ion se obtiene
\[
 \mathrm{ENC}=\frac{\sigma_D}{g},\qquad
 \frac{S}{N}=\frac{Ng}{\sigma_D}=\frac{N}{\mathrm{ENC}}.
\]
El ENC se expresa en iones de carga $+e$ RMS.
Se calcula el umbral para S/N=5 y se comprueba con transitorios LTspice
de números enteros de iones a ambos lados de ese umbral.

\section{Resultados}

Con promedios de 300 µs se obtiene un ENC de 40.47 iones RMS.
El umbral para $S/N>5$ es de 203 iones: la simulación da $S/N=4.989$
para 202 iones y $S/N=5.014$ para 203.

También se ha comprobado la contribución de una población igual de iones
desplazada un frame. La población del frame anterior
contribuye un 2.88\% de la señal del frame de interés; la del siguiente,
un 1.01\%. Estas contribuciones incluyen la inducción antes de la llegada,
la memoria del circuito y el promedio temporal.

\section{Comparación del ruido flicker}

Se ha repetido el cálculo del PDF  \texttt{Noise Analysis CSA 180 nm.pdf}
con el esquema original: \texttt{Cd1=250 fF}, temperatura por defecto de
27 °C y barrido de 1 Hz a 1 GHz con 1000 puntos por década.

{
\begin{longtable}[]{@{}
  >{\raggedleft\arraybackslash}p{(\linewidth - 4\tabcolsep) * \real{0.3333}}
  >{\raggedleft\arraybackslash}p{(\linewidth - 4\tabcolsep) * \real{0.3333}}
  >{\raggedleft\arraybackslash}p{(\linewidth - 4\tabcolsep) * \real{0.3333}}@{}}
\toprule\noalign{}
\begin{minipage}[b]{\linewidth}\raggedleft
Frecuencia
\end{minipage} & \begin{minipage}[b]{\linewidth}\raggedleft
Ruido de salida reproducido
\end{minipage} & \begin{minipage}[b]{\linewidth}\raggedleft
Ruido equivalente de corriente de entrada
\end{minipage} \\
\midrule\noalign{}
\endhead
\bottomrule\noalign{}
\endlastfoot
1 Hz & 2.167 \(\mathrm{mV}/\sqrt{\mathrm{Hz}}\) & 0.09287
\(\mathrm{fA}/\sqrt{\mathrm{Hz}}\) \\
10 Hz & 0.2187 \(\mathrm{mV}/\sqrt{\mathrm{Hz}}\) & 0.09288
\(\mathrm{fA}/\sqrt{\mathrm{Hz}}\) \\
100 Hz & 0.02207 \(\mathrm{mV}/\sqrt{\mathrm{Hz}}\) & 0.09373
\(\mathrm{fA}/\sqrt{\mathrm{Hz}}\) \\
\end{longtable}
}

A 1 Hz se obtiene $2.167\,\mathrm{mV}/\sqrt{\mathrm{Hz}}$, compatible con
los aproximadamente $2.2\,\mathrm{mV}/\sqrt{\mathrm{Hz}}$ de la figura que mandas.
En la ejecución del estudio, las siete contribuciones \texttt{1overf} son cero.
El fichero de modelos recibido deja sin especificar \texttt{KF},
\texttt{NOIA}, \texttt{NOIB}, \texttt{NOIC} y \texttt{NOIMOD}.

Se comprueba la lectura de esos términos con una copia de los modelos y un
valor de prueba \texttt{KF=1e-25}: las siete contribuciones pasan a ser
no nulas. En ambas ejecuciones, la suma de las varianzas por transistor
reproduce la varianza total.

Entre 1 y 10 Hz, la corriente de ruido equivalente es casi constante y el
ruido de salida cae un factor diez por la respuesta del circuito. Para
identificar la contribución flicker es necesario comparar los términos
\texttt{1overf} de ambas ejecuciones. Para ello se necesitan la versión de LTspice,
el netlist, los modelos y los datos exportados de esas contribuciones y
del espectro total de la ejecución.

%\section{Cálculos adjuntos}
%
%El script \path{scripts/build_capacitance_study.py} reproduce el cálculo del ENC;
%\path{scripts/audit_ltspice_noise.py} reproduce la comparación de ruido.
%Los resultados, espectros y netlists están en \path{results/ltspice/}.
%Las comprobaciones de carga integrada, paso temporal e integración espectral
%están guardadas junto a los resultados en los archivos JSON.

\end{document}
